\SourceLeaf{107}tum magnetismi liberi primae acus $\Sigma ae$ per $m$ denotamus, quum constet esse $\Sigma be=0$: pars ipsius $\frac{\mathrm{d}\Omega}{\mathrm{d}u}$ e termino primo ipsius $\Omega$ redundans erit $=-mT\sin u$.

Statuendo brevitatis caussa:
\[
\begin{aligned}
k={}&\alpha\cos\psi+\beta\sin\psi+A\cos(\psi-U)+B\sin(\psi-U)\\
&-a\cos(\psi-u)-b\sin(\psi-u)\\
l={}&(\alpha\sin\psi-\beta\cos\psi+A\sin(\psi-U)-B\cos(\psi-U)\\
&-a\sin(\psi-u)+b\cos(\psi-u))^2+(C-c)^2
\end{aligned}
\]
erit $rr=(R+k)^2+l$.

Quum in experimentis utilibus $R$ dimensionibus utriusque acus multo maior esse debeat, quantitas $\frac1{r^{n-1}}$ in seriem valde convergentem
\[
\begin{aligned}
&R^{-(n-1)}-(n-1)kR^{-n}+\left(\frac{nn-n}{2}kk-\frac{n-1}{2}l\right)R^{-(n+1)}\\
&\qquad-\left(\tfrac16(n^3-n)k^3-\tfrac12(nn-1)kl\right)R^{-(n+2)}+\text{etc.}
\end{aligned}
\]
evolvitur, cuius lex, si operae pretium esset, facile assignaretur. Singuli termini aggregati $\Sigma\frac{eE}{r^{n-1}}$, post substitutionem valorum quantitatum $k,l$ prodeuntes implicabunt factorem talem
\[
\Sigma eEa^\lambda b^\mu c^\nu A^{\lambda'}B^{\mu'}C^{\nu'}
\]
qui aequivalet producto e factoribus $\Sigma ea^\lambda b^\mu c^\nu$, $\Sigma EA^{\lambda'}B^{\mu'}C^{\nu'}$ a statu magnetico primae et secundae acus resp. pendentibus. Quae hoc respectu generaliter stabilire licet, restringuntur ad aequationes
\[
\begin{gathered}
\Sigma e=0,\ \Sigma ea=m,\ \Sigma eb=0,\ \Sigma ec=0,\\
\Sigma E=0,\ \Sigma EA=M,\ \Sigma EB=0,\ \Sigma EC=0
\end{gathered}
\]
ubi per $M$ denotamus momentum magnetismi liberi secundae acus. In casu speciali, ubi acus prioris figura magnetismique distributio est symmetrica iuxta longitudinem, puta ut bina semper elementa sibi respondeant, pro quibus $a$ et $e$ habeant valores oppositos, $b$ et $c$ aequales, centro cum puncto $h$ coincidente, semper erit $\Sigma ea^\lambda b^\mu c^\nu=0$ pro valore pari numeri $\lambda+\mu+\nu$, et similia valent de secunda acu, si figura magnetismique distributio respectu puncti $H$ symmetrica est. Generaliter itaque evanescent in aggregato $\Sigma\frac{eE}{r^{(n-1)}}$ coëfficientes potestatum $R^{-(n-1)}$ et $R^{-n}$; in casu speciali, ubi utraque acus symmetrica magnetismoque symmetrice imbuta est, simulque centrum prioris, $h$ et $h'$, nec non centrum posterioris et $H$ coincidant, evanescent etiam coëfficientes potestatum $R^{-(n+2)},R^{-(n+4)},R^{-(n+6)}$ etc., qui, quoties conditiones illae proxime locum ha\SourceLeaf{108}bent, saltem perparvi evadere debent. Terminus principalis, qui ex evolutione partis secundae ipsius $\Omega$, puta huius $-\Sigma\frac{eE}{(n-1)r^{(n-1)}}$, prodit, erit
\[
\begin{aligned}
&=-\tfrac12 R^{-(n+1)}(n\Sigma eEkk-\Sigma eEl)\\
&=mMR^{-(n+1)}(n\cos(\psi-U)\cos(\psi-u)-\sin(\psi-U)\sin(\psi-u))
\end{aligned}
\]

Hinc colligitur, partem ipsius $\frac{\mathrm{d}\Omega}{\mathrm{d}u}$ actioni acus secundae respondentem exprimi per seriem talem
\[
fR^{-(n+1)}+f'R^{-(n+2)}+f''R^{-(n+3)}+\text{etc.}
\]
ubi coëfficientes sunt functiones rationales cosinuum et sinuum angulorum $\psi,u,U$ atque quantitatum $\alpha,\beta$, insuperque implicant quantitates constantes a statu magnetico acuum pendentes; et quidem erit
\[
f=mM(n\cos(\psi-U)\sin(\psi-u)+\sin(\psi-U)\cos(\psi-u))
\]
Evolutio completa coëfficientium sequentium $f',f''$ etc. ad institutum nostrum non est necessaria: sufficit observare

1) in casu symmetriae perfectae modo addigitatae coëfficientes $f',f'''$ etc. evanescere,

2) si manentibus quantitatibus reliquis invariatis $\psi$ augeatur duobus rectis (sive quod idem est, si distantia $R$ capiatur in eadem recta retrorsum producta ab altera parte puncti $h'$), coëfficientes $f,f'',f''''$ etc. valores suos retinere, contra $f',f''',f^{\mathrm{v}}$ etc. valores oppositos nancisci, sive seriem in
\[
fR^{-(n+1)}-f'R^{-(n+2)}+f''R^{-(n+3)}-\text{etc.}
\]
mutari: facile hoc inde concluditur, quod per illam mutationem ipsius $\psi$, $k$ transit in $-k$, $l$ vero non mutatur.

\Article{17}
Conditio itaque, ut acus mobilis per complexum virium non vertatur circa axem verticalem, comprehenditur in aequatione sequente
\[
0=-mT\sin u+fR^{-(n+1)}+f'R^{-(n+2)}+f''R^{-(n+3)}+\text{etc.}-\theta(u-N)
\]

Quum facile effici possit, ut valor ipsius $N$, si non exacte $=0$, saltem perparvus sit, atque etiam $u$ pro experimentis, de quibus hic agitur, intra arctos \SourceLeaf{109}limites maneat, pro termino $\theta(u-N)$ absque erroris sensibilis metu substituere licebit $\theta\sin(u-N)$, eo magis, quod $\frac{\theta}{mT}$ est fractio perparva. Sit $u^0$ valor ipsius $u$, aequilibrio acus primae absente secunda respondens, sive
\[
mT\sin u^0+\theta\sin(u^0-N)=0
\]
unde facile colligitur
\[
mT\sin u+\theta\sin(u-N)=(mT\cos u^0+\theta\cos(u^0-N))\sin(u-u^0)
\]
ubi loco factoris primi tuto adoptare licet $mT+\theta$. Ita aequatio nostra fit
\[
(mT+\theta)\sin(u-u^0)=fR^{-(n+1)}+f'R^{-(n+2)}+f''R^{-(n+3)}+\text{etc.}
\]
Quodsi hic terminum primum $fR^{-(n+1)}$ solum retinemus, solutio in promptu est, scilicet habemus
\[
\begin{aligned}
&\operatorname{tang}(u-u^0)\\
&=\frac{\begin{gathered}mM(n\cos(\psi-U)\sin(\psi-u^0)\\+\sin(\psi-U)\cos(\psi-u^0))R^{-(n+1)}\end{gathered}}{\begin{gathered}mT+\theta+mM(n\cos(\psi-U)\cos(\psi-u^0)\\-\sin(\psi-U)\sin(\psi-u^0))R^{-(n+1)}\end{gathered}}
\end{aligned}
\]
ubi in denominatore partem, quae implicat factorem $R^{-(n+1)}$, eodem iure supprimere poterimus, sive statuere
\[
\begin{aligned}
\operatorname{tang}(u-u^0)&=\frac{mM}{mT+\theta}(n\cos(\psi-U)\sin(\psi-u^0)\\
&\qquad+\sin(\psi-U)\cos(\psi-u^0))R^{-(n+1)}\\
&=FR^{-(n+1)}
\end{aligned}
\]

Si vero terminos ulteriores respicere volumus, patet, $\operatorname{tang}(u-u^0)$ in seriem talem evolvi
\[
\operatorname{tang}(u-u^0)=FR^{-(n+1)}+F'R^{-(n+2)}+F''R^{-(n+3)}+\text{etc.}
\]
ubi levis attentio docet, coëfficientes $F,F',F''$ etc. usque ad coëfficientem potestatis $R^{-(2n+1)}$ incl. oriri resp. ex
\[
\frac{f}{mT+\theta},\quad\frac{f'}{mT+\theta},\quad\frac{f''}{mT+\theta}\quad\text{etc.}
\]
mutato $u$ in $u^0$, inde a termino sequente autem partes novae accedent, quibus tamen accuratius persequendis ad institutum nostrum non opus est. Ceterum manifesto $u-u^0$ in seriem similis formae explicabitur, quae adeo usque ad potestatem $R^{-(3n+2}$ cum serie pro $\operatorname{tang}(u-u^0)$ coincidet.

\SourceLeaf{110}
\Article{18}
Patet iam, si acu secunda in diversis punctis eiusdem rectae collocata, ut manentibus $\psi$ et $U$ sola distantia $R$ mutetur, deflexiones acus mobilis a statu aequilibrii, acu secunda absente, puta anguli $u-u^0$ observentur, hinc valores coëfficientium $F,F',F''$ etc., quotquot adhuc sensibiles sunt, per eliminationem erui posse, quo facto habebimus
\[
\frac MT=\left(1+\frac{\theta}{Tm}\right)\frac{F}{n\cos(\psi-U)\sin(\psi-u^0)+\sin(\psi-U)\cos(\psi-u^0)}
\]
ubi valor quantitatis $\frac{\theta}{Tm}$ per methodum, quam in art. $8$ docuimus, inveniri poterit. Sed ad praxin magis commodam sequentia observare e re erit.

I. Loco comparationis ipsius $u$ cum $u^0$ praestat, binas deflexiones oppositas inter se comparare, situ acus secundae inverso, puta, ut manentibus $R$ et $\psi$ angulus $U$ duobus rectis augeatur. Designatis valoribus ipsius $u$ his positionibus respondentibus per $u',u''$, \textit{exacte} foret $u''=-u'$ pro casu symmetriae perfectae, si simul $u^0=0$; sed superfluum est, has conditiones anxie servare, quum pateat, $u'$ et $u''$ per series similes determinari, in quibus termini primi \textit{exacte} oppositos valores habeant, adeoque etiam $\tfrac12(u'-u'')$ nec non $\operatorname{tang}\tfrac12(u'-u'')$ per seriem similem, in qua termini primi coëfficiens sit \textit{exacte} $=F$.

II. Adhuc melius erit, quaterna semper experimenta copulare, etiam angulo $\psi$ duobus rectis mutato, sive distantia $R$ ab altera parte sumta. Si duobus posterioribus experimentis respondent valores ipsius $u$ hi $u''',u''''$, etiam differentia $\tfrac12(u'''-u'''')$ per similem seriem exprimetur, cuius terminus primus quoque habebit coëfficientem $=F$. Observare convenit (quod e praecedentibus facile colligitur), si $n$ esset numerus impar, coëfficientes $F,F'',F''''$ etc. in infinitum in utraque serie pro $u'-u^0$, atque $u'''-u^0$ exacte aequales, coëfficientesque $F',F''',F^{\mathrm{v}}$ etc. in infinitum exacte oppositos fore, et perinde pro $u''-u^0$ et $u''''-u^0$, ita ut in serie pro $u'-u''+u'''-u''''$ termini alternantes exciderent. Sed in casu naturae, ubi $n=2$, generaliter loquendo ista relatio inter series pro $u'-u^0$ atque $u'''-u^0$ stricte non valet, quum iam pro potestate $R^{-6}$ coëfficientes non exacte oppositi prodeant: attamen ostendi potest, pro hoc quoque termino compensationem completam in combinatione $u'-u''+u'''-u''''$ intercedere, ita ut $\operatorname{tang}\tfrac14(u'-u''+u'''-u'''')$ habeat formam
\[
LR^{-3}+L'R^{-5}+L''R^{-7}+\text{etc.}
\]
\SourceLeaf{111}sive generalius, dum valorem ipsius $n$ tantisper indeterminatum linquimus, hancce
\[
LR^{-(n+1)}+L'R^{-(n+3)}+L''R^{-(n+5)}+\text{etc.}
\]
existente $L=F$.

III. Angulos $\psi,U$ ita eligere expedit, ut leves errores in ipsorum mensura commissi valorem ipsius $F$ sensibiliter non mutent. Ad hunc finem valor ipsius $U$, pro valore \textit{dato} ipsius $\psi$, ita accipi debet, ut $F$ fiat maximum, puta esse debet
\[
\operatorname{cotang}(\psi-U)=n\operatorname{tang}(\psi-u^0)
\]
unde fit
\[
F=\pm\frac{mM}{mT+\theta}\BareRoot{(nn\sin(\psi-u^0)^2+\cos(\psi-u^0)^2)}
\]

Angulus vero $\psi$ ita eligendus est, ut hic valor ipsius $F$ fiat vel maximum vel minimum: illud evenit pro $\psi-u^0=90^0$ vel $270^0$, ubi
\[
F=\pm\frac{nmT}{mT+\theta}
\]
hoc pro $\psi-u^0=0$ vel $180^0$, ubi
\[
F=\pm\frac{mM}{mT+\theta}
\]

\Needspace{18\baselineskip}
\Article{19}
Duae itaque methodi praesto sunt ad praxin maxime idoneae, quarum elementa sequens schema exhibet.

\expandafter\def\csname HistoricalNoteText11\endcsname{Accuratius, ad planum verticale, cui respondet valor $u=u^0$, i. e. in quo axis magneticus in aequilibrio est, acu secunda absente. Ceterum in praxi, differentia tum propter parvitatem, tum propter ipsam rationem, a qua in art. praec. III. profecti sumus, tuto semper negligi potest.}
\begin{center}
\textit{Modus primus.}

Acus secundae tum centrum tum axis in recta ad meridianum magneticum\HistoricalNote{11} normali.

\begin{tabular}{c|c|c|l|l}
& & &centrum&polus borealis\\
Deflexio&\multicolumn{2}{c|}{Situs acus}&versus&versus\\\hline
$u=u'$&$\psi=90^0$&$U=90^0$&orientem&orientem\\
$u=u''$&$\psi=90^0$&$U=270^0$&orientem&occidentem\\
$u=u'''$&$\psi=270^0$&$U=90^0$&occidentem&orientem\\
$u=u''''$&$\psi=270^0$&$U=270^0$&occidentem&occidentem
\end{tabular}
\end{center}
